\documentclass[11pt]{article}
\usepackage[margin=1.1in]{geometry}
\usepackage{amsmath,amssymb,amsthm}
\usepackage{natbib}
\usepackage{array}
\usepackage{booktabs}
\usepackage{xcolor}
\usepackage[colorlinks=true,linkcolor=blue!50!black,citecolor=blue!50!black,urlcolor=blue!50!black]{hyperref}

\makeatletter
\renewcommand\paragraph{\@startsection{paragraph}{4}{\z@}%
  {3.25ex \@plus1ex \@minus.2ex}{-1em}%
  {\normalfont\normalsize\itshape}}
\makeatother

\newtheorem{proposition}{Proposition}
\newtheorem{corollary}{Corollary}
\theoremstyle{definition}
\newtheorem{definition}{Definition}
\theoremstyle{remark}
\newtheorem{remark}{Remark}

\newcommand{\E}{\mathbb{E}}

\title{Consistent Objectives for Self-Improving Systems \\[4pt]
\large Absolute, Relative and Evolutionary Assessment Over Short and Long Horizons}
\author{Peter Cotton\thanks{Working draft, October 2026. Numerical certificates:
\texttt{papers/objectives/verify\_objectives.py} in
\url{https://github.com/microprediction/metareasoning}.}}
\date{}

\begin{document}
\maketitle

\begin{abstract}
A self-improving system accepts or rejects changes to itself by a per-step objective, and is
judged by where it ends up. We ask which per-step objectives select the same trajectories as
that long-term goal, comparing absolute scores, win rates against a predecessor or a fixed
anchor, selection of a population's best, and relative wealth under compounding. Improvement
in a fixed absolute score, less a constant cost per step, is the only per-step reward
consistent for every feasible set of trajectories. Compounded relative wealth is consistent
with expected log return, and the other relative schemes reward churn, distort attitudes to
risk, reward variance, or record evaluator drift as progress.
\end{abstract}

\section{Introduction}

Each proposed change to a self-improving system, whether to a prompt, a memory, a piece of
code or a set of weights, is accepted or rejected on evidence available now.
The system as a whole is judged later, by what it has become. The two judgements coincide only
for particular choices of the per-step objective.

Some objectives in use are \emph{absolute}, such as a benchmark score, a test pass rate, or
a proper scoring rule against resolved outcomes. Others are \emph{relative}. Self-play and
self-rewarding language models score a version against its predecessor
\citep{yuan2024selfrewarding}. Leaderboards rank versions by pairwise preference
\citep{chiang2024arena}. Evolutionary search over agents keeps the survivors of a population
\citep{zhang2025dgm,iacob2026redqueen}.

The literature on each line is large. Reward shaping characterises which per-step rewards
preserve optimal policies \citep{ng1999shaping}. Goodhart's law, reward hacking and reward-model
overoptimisation describe what goes wrong when a proxy is optimised hard
\citep{manheim2018goodhart,skalse2022hacking,gao2022overoptimization,karwowski2023goodhart}.
Non-transitive games show that beating a predecessor need not mean getting stronger
\citep{balduzzi2019openended,czarnecki2020spinning}. Forecasting competitions break the
propriety of scoring rules \citep{lichtendahl2007competition,witkowski2021competitions}.
Evolutionary finance shows which beliefs survive market selection
\citep{blume1992evolution,sandroni2000markets,blume2006smart}.

The literature rarely puts these results side by side. Recent frameworks classify self-improving systems
by whether their standard is grounded outside the agent \citep{tang2026gai} and catalogue
evaluator drift as a safety risk \citep{gong2026evolutionary}. The Red Queen G\"odel Machine
holds the evaluator fixed within epochs so that improvement guarantees hold per epoch
\citep{iacob2026redqueen}. None of them compares the objectives themselves for consistency
between the step and the horizon. We make that comparison with elementary propositions, each
checked numerically.

\section{Setup}

A system passes through versions $\theta_0,\theta_1,\dots,\theta_T$. At each step it chooses
$\theta_{t+1}$ from a feasible set $N(\theta_t)$, which may or may not include staying put. The \emph{long-term
objective} is a utility of the final version, $U(\theta_T)$. For a discounted sum of deployed
utility, the potential term in Proposition~\ref{prop:potential} becomes
$\gamma\Phi(\theta')-\Phi(\theta)$ \citep{ng1999shaping}. A \emph{per-step reward}
$R(\theta,\theta')$ scores a single change, and the system maximises
$\sum_{t<T} R(\theta_t,\theta_{t+1})$, or a greedy approximation to it.

\begin{definition}[Horizon consistency]
A per-step reward $R$ is \emph{horizon consistent} with $U$ if, for every horizon $T$, every
starting version and every feasibility structure $N$, the trajectories that maximise
$\sum_{t<T}R(\theta_t,\theta_{t+1})$ are exactly those that maximise $U(\theta_T)$.
\end{definition}

The definition asks for agreement on the whole set of maximisers. A reward that agrees with $U$
on most problems but not all fails on the problems a thorough search finds, which is the
mechanism behind Goodhart's law \citep{manheim2018goodhart}.

\section{Absolute assessment}

\begin{proposition}[Only potential differences are consistent]\label{prop:potential}
A per-step reward $R$ is horizon consistent with $U$ if and only if
\[
R(\theta,\theta') = \Phi(\theta') - \Phi(\theta) + k
\]
for some constant $k$ and some strictly increasing transformation $\Phi$ of $U$.
\end{proposition}

The reward for a change is the rise in a monotone transformation of the long-term score, plus
the same constant for every step.

\begin{proof}
Sufficiency is telescoping: the sum over a path of length $T$ equals
$\Phi(\theta_T)-\Phi(\theta_0)+kT$, which ranks endpoints as $U$ does.

For necessity, take $T=1$
and let $N(\theta_0)$ be an arbitrary set of versions: $R(\theta_0,\cdot)$ must rank versions as
$U$ does. Take $T=2$ with $N(a)=\{b,b'\}$ and $N(b)=N(b')=\{c\}$. The two feasible paths
$a\to b\to c$ and $a\to b'\to c$ share their endpoint, so both maximise $U$ and they must tie, and $R(a,b)+R(b,c)$ cannot depend on
$b$.

Write it as $h(a,c)$. Fixing $c=c_0$ gives $R(a,b)=\alpha(a)-\beta(b)$ with
$\alpha(a)=h(a,c_0)$ and $\beta(b)=R(b,c_0)$. Substituting back, $\alpha(b)-\beta(b)$ must be a
constant $k$, so $R(a,b)=\alpha(a)-\alpha(b)+k$. Set $\Phi=-\alpha$. The $T=1$ condition makes
$\Phi$ increasing in $U$.
\end{proof}

A negative $k$ is a cost per step, and $-k$ can be read as the price of a computation in the
sense of \citet{russell1991metareasoning}. Proposition~\ref{prop:potential} is the undiscounted
case of the shaping theorem of \citet{ng1999shaping}, specialised to a system whose state is its
own version. Under uncertainty only affine $\Phi$ preserve rankings of expected utility, so in
practice the consistent reward is improvement in the long-term score itself, less a cost.

\paragraph{The score must be proper.} If $U$ is an expected score of probabilistic forecasts
against resolved outcomes, a strictly proper scoring rule makes truthful reporting optimal
\citep{brier1950verification,gneiting2007proper}. An improper score adds a second
inconsistency, between the forecasts the system holds and the forecasts it reports. Training on
resolved outcomes with a proper score \citep{turtel2026futureaslabel} meets both conditions.

\paragraph{Greedy acceptance is a separate failure.} Proposition~\ref{prop:potential} concerns
the objective, not the search. A system that accepts the change with the largest immediate
improvement can still be arbitrarily far from optimal. If one change gains $\epsilon$ and leads
nowhere while another gains nothing and unlocks a gain of one, greedy acceptance takes the
first. The remedy is to value a change by its effect on what can follow, which is the value of
computation \citep{russell1991metareasoning,hay2012selecting}. For changes that run in
parallel, the corresponding rule keeps a path while its future pivotal value exceeds its
carrying cost \citep{cotton2026wavefront}.

\section{Relative assessment against a predecessor}

Pairwise comparison is attractive because a judge can often say which of two outputs is better
when it cannot score either. Suppose comparisons follow a Bradley--Terry model
\citep{bradley1952rank}, so that $\theta'$ beats $\theta$ with probability
$\sigma(u(\theta')-u(\theta))$ for a latent strength $u$ and the logistic function $\sigma$.

\begin{proposition}[Win rate against the predecessor rewards churn]\label{prop:churn}
Let $R(\theta,\theta')=\sigma(u(\theta')-u(\theta))$.
\begin{enumerate}
\item Accepting a change exactly when $R>\tfrac12$ accepts exactly the changes that raise $u$.
\item Splitting any gain $D>0$ into two steps of $D/2$ raises the cumulative reward from
$\sigma(D)<1$ to $2\sigma(D/2)>1$. Two changes that leave $u$ unchanged earn $1$, more than any
single change of any size.
\item $R$ is not horizon consistent. There are feasibility structures in which the cumulative
reward selects an endpoint of lower $u$.
\end{enumerate}
\end{proposition}

\begin{proof}
The first part is monotonicity of $\sigma$. The second uses $\sigma(0)=\tfrac12$ and
$\sigma<1$. For the third, let the start have $u=0$ and offer two routes over four steps: a jump
to $u=5$ followed by staying put, or a ladder of four unit steps to $u=4$. The jump route earns
$\sigma(5)+\tfrac32\approx 2.49$ and the ladder $4\sigma(1)\approx 2.92$. The cumulative win
rate takes the ladder to the worse endpoint.
\end{proof}

The ordinal acceptance rule is sound and the cumulative objective is not. A system trained to
maximise accumulated wins against its own predecessors is paid for the number of changes, not
for their size. Without a transitive latent strength the failure is worse: in a game with
strategic cycles, every step can beat its predecessor while the sequence returns to its start
\citep{balduzzi2019openended}. Real games combine a transitive axis with many cycles at each
level of strength \citep{czarnecki2020spinning}, which is why population methods evaluate
against mixtures rather than the last version
\citep{lanctot2017unified,balduzzi2018reevaluating}.

\section{Relative assessment against a fixed anchor}

Comparing every version with one fixed reference, $R=\sigma(u(\theta)-u_{\mathrm{ref}})$, is a
strictly increasing transformation of $u$. It is horizon consistent when outcomes are certain.
It is not affine, so it changes the system's attitude to risk.

\begin{proposition}[A fixed anchor makes laggards risk-seeking]\label{prop:anchor}
If the uncertain strength of a proposed version has support entirely below $u_{\mathrm{ref}}$,
any mean-preserving spread raises the expected win rate against the anchor. If the support lies
entirely above, any mean-preserving spread lowers it.
\end{proposition}

\begin{proof}
The logistic function $\sigma$ is convex on $(-\infty,0)$ and concave on $(0,\infty)$, and
Jensen's inequality gives both claims.
\end{proof}

In a normal example with the anchor three units ahead, nearly all of the mass lies where
$\sigma$ is convex. Widening the standard deviation of a proposal from $0.1$ to $1$ raises its
expected win rate from $0.048$ to $0.069$. Three units behind, the same spread
lowers it from $0.952$ to $0.931$. A system that trails its anchor is paid to gamble, and a
system that leads is paid to stop experimenting.

\section{Relative assessment within a population}

Evolutionary and tournament methods keep the best of several candidates, judged on realised
scores. The candidate with the highest expected score is not the one most likely to post the
highest realised score.

\begin{proposition}[Selection of the realised best rewards variance]\label{prop:variance}
Let candidate $i$ post $S_i=\mu_i+s_iZ_i$ with independent standard normal $Z_i$. For two
candidates,
\[
P(S_1>S_2)=\mathcal{N}\!\left(\frac{\mu_1-\mu_2}{\sqrt{s_1^2+s_2^2}}\right),
\]
where $\mathcal{N}$ is the standard normal distribution function. The probability increases
in $s_1$ whenever $\mu_1<\mu_2$.
\end{proposition}

\begin{proof}
The difference $S_1-S_2$ is normal with mean $\mu_1-\mu_2$ and variance $s_1^2+s_2^2$, which
gives the formula. When $\mu_1<\mu_2$ the argument of $\mathcal{N}$ is negative, and raising
$s_1$ moves it towards zero.
\end{proof}

A trailing candidate gains by being noisier. With $\mu_2-\mu_1=0.5$ and $s_2=1$, raising $s_1$
from $0.5$ to $4$ lifts its chance of selection from $0.33$ to $0.45$. Take a field of ten
candidates with unit standard deviation and equal means, except that one is better by half a
unit. Doubling the standard deviation of another candidate makes it the most likely winner: it
is selected $23\%$ of the time, against $17\%$ for the better candidate. Closed forms and Monte Carlo
agree to within $2\times10^{-3}$ in win probability.

Forecasters competing for a single prize report more extreme probabilities than they believe
\citep{lichtendahl2007competition}, for the same reason, and incentive-compatible competitions
are designed to remove the distortion \citep{witkowski2021competitions}. The selected
alternative's realised score overstates its value \citep{smith2006optimizer}. In backtests the
best of many strategies needs a deflated threshold \citep{bailey2014deflated}, and LLM-discovered
trading strategies fail once the agent's own trial count is accounted for
\citep{gencay2026honest}.

\paragraph{When rewarding variance is right.} If the long-term objective is itself an expected
maximum, selection by realised maximum is consistent with it. Best-of-$k$ deployment and
pass@$k$ pay the best of several attempts, and there variance has value
\citep{cotton2026wavefront,cotton2026passk}. Selection by maximum suits a deployment that pays the
maximum, and misleads one that pays the mean.

\paragraph{Finite populations and spite.} Relative fitness in a finite population rewards
lowering a rival's payoff as much as raising one's own. \citet{hamilton1970spite} identified
spiteful behaviour, and \citet{schaffer1988finite} showed that evolutionarily stable strategies
in finite populations maximise relative rather than absolute payoff. A population of agents
judged against one another can therefore select changes that make every member worse in
absolute terms.

\section{Relative assessment under compounding}

Relative assessment becomes consistent with an absolute objective when wealth compounds. Let
strategies compound
wealth by gross returns $R_{i,t}$, independent over time with finite $\E|\log R_i|$, and let
assessment be relative: the share $w_{i,t}/\sum_j w_{j,t}$.

\begin{proposition}[Relative wealth selects the largest expected log return]\label{prop:kelly}
If $\E\log R_1>\E\log R_j$ for all $j\neq1$, the share of strategy 1 tends to one almost
surely, whatever the expected returns $\E R_j$.
\end{proposition}

\begin{proof}
By the strong law of large numbers, $t^{-1}\log w_{i,t}\to\E\log R_i$ almost surely. The ratio
$w_{j,t}/w_{1,t}$ therefore tends to zero for every $j\neq1$.
\end{proof}

The proposition is the selection side of \citet{kelly1956information} and
\citet{breiman1961optimal}. It separates the short-term and long-term readings of an absolute
objective.

A strategy with returns $1.6$ or $0.5$ with equal probability has $\E R=1.05$ and
$\E\log R=-0.11$. A strategy with returns $1.1$ or $0.95$ has $\E R=1.025$ and
$\E\log R=+0.022$. The first wins on one-step expected return, and its expected wealth after
400 steps exceeds $10^8$. It loses the population in every one of $20{,}000$ simulated runs,
since its median log wealth falls by $0.11$ a step. Expected wealth and typical wealth
diverge because wealth is not ergodic \citep{peters2019ergodicity}.

\paragraph{Beliefs.} In markets the same selection operates on forecasts. Agents survive in
proportion to how well their beliefs track the truth
\citep{blume1992evolution,sandroni2000markets}, with qualifications for incomplete markets
\citep{blume2006smart}. A forecaster paid by the log score relative to a market grows its share
at the rate of its log-score advantage, so compounded relative assessment recovers the log
score, a strictly proper absolute objective. \citet[ch.~5]{cotton2022microprediction} considers agents that solicit, compensate and
combine the predictions of other agents. When such compensation is relative and compounds, it
belongs to this class.

\paragraph{Absolute selection.} When fitness does not depend on the population, selection
raises mean fitness. Under replicator dynamics \citep{taylor1978ess} with constant
fitnesses, mean fitness rises at a rate equal to the variance of fitness. This is the simplest
case of Fisher's fundamental theorem, whose precise statement is due to
\citet{price1972fisher}. The guarantee does not survive frequency-dependent fitness, the setting
of evolutionary game theory \citep{maynardsmith1973conflict}.

\section{Moving evaluators}

A system that helps write its own evaluator breaks the telescoping in
Proposition~\ref{prop:potential}. Let $\Phi_t$ be the evaluator in force at step $t$, and let the
system record the improvement $\Phi_t(\theta_{t+1})-\Phi_t(\theta_t)$ at each step.

\begin{proposition}[Measured progress includes evaluator drift]\label{prop:drift}
\[
\sum_{t<T}\big[\Phi_t(\theta_{t+1})-\Phi_t(\theta_t)\big]
=\Phi_T(\theta_T)-\Phi_0(\theta_0)-\sum_{t<T}\big[\Phi_{t+1}(\theta_{t+1})-\Phi_t(\theta_{t+1})\big].
\]
\end{proposition}

The identity follows by adding and subtracting $\Phi_{t+1}(\theta_{t+1})$ at each step. Recorded
progress equals the change in the current evaluator's opinion, less the drift in the evaluator
at the versions actually visited. With a fixed anchored evaluator the drift term vanishes and
recorded progress is real.

The drift term has either sign. An evaluator that tightens as the system improves makes recorded
progress understate the change, which is the intent of co-evolving evaluators
\citep{iacob2026redqueen}. An evaluator that comes to share the system's errors, as when a model
judges its own outputs \citep{yuan2024selfrewarding,huang2023selfcorrect} or trains on them
\citep{shumailov2024collapse}, can make the anchored score stand still while the measured one
rises. Holding the evaluator fixed within an epoch restores telescoping within the epoch, and
accumulating real data rather than replacing it avoids collapse in the settings studied by
\citet{gerstgrasser2024collapse}.

\section{Objectives compared}

\begin{center}
\small
\renewcommand{\arraystretch}{1.25}
\begin{tabular}{@{}>{\raggedright\arraybackslash}p{0.33\textwidth}>{\raggedright\arraybackslash}p{0.28\textwidth}>{\raggedright\arraybackslash}p{0.3\textwidth}@{}}
\toprule
Per-step objective & Consistent with the long-term goal & Failure when optimised \\
\midrule
Improvement in a fixed proper score, less cost & Yes, for every feasibility structure & Greedy search misses option value \\
Win rate against the predecessor & Ordinal acceptance only & Churn, and cycles if non-transitive \\
Win rate against a fixed anchor & Yes under certainty & Risk-seeking behind, risk-averse ahead \\
Best of a population on realised score & Only for expected-maximum goals & Rewards variance, winner's curse \\
Relative fitness in a finite population & No & Spite \\
Relative wealth under compounding & Yes, for expected log return & One-step expected value loses \\
Improvement under a moving evaluator & Only with a fixed anchor & Drift recorded as progress \\
\bottomrule
\end{tabular}
\end{center}

\section{Closing}

Neither absolute nor relative assessment dominates. Absolute improvement in a fixed proper score
is the one per-step objective consistent for every feasibility structure, but it needs a score
that can be computed, and a search that values option value as well as immediate gain.
Relative assessment is cheaper and often the only signal available. It is consistent in two
cases, ordinal acceptance under a transitive latent strength and multiplicative compounding.
Elsewhere it produces churn, distorted risk, rewarded variance, spite or drift.

The trade-off is not particular to machine learning. Rank-order tournaments for workers
\citep{lazear1981tournaments}, forecasting contests and natural selection face the same
choice between paying for level and paying for rank. We leave to future work the empirical
measurement of these effects in a running self-improving system, where the drift term in
Proposition~\ref{prop:drift} can be estimated by re-scoring past versions under the current
evaluator.

\begin{thebibliography}{}\small

\bibitem[Bailey and L\'opez de Prado(2014)]{bailey2014deflated}
D.~H. Bailey and M.~L\'opez de Prado. The deflated Sharpe ratio: correcting for selection bias,
backtest overfitting and non-normality. \emph{Journal of Portfolio Management}, 40(5):94--107,
2014. \href{https://doi.org/10.3905/jpm.2014.40.5.094}{doi:10.3905/jpm.2014.40.5.094}.

\bibitem[Balduzzi et~al.(2018)]{balduzzi2018reevaluating}
D.~Balduzzi, K.~Tuyls, J.~Perolat and T.~Graepel. Re-evaluating evaluation. In \emph{NeurIPS},
2018. \href{https://arxiv.org/abs/1806.02643}{arXiv:1806.02643}.

\bibitem[Balduzzi et~al.(2019)]{balduzzi2019openended}
D.~Balduzzi, M.~Garnelo, Y.~Bachrach, W.~M. Czarnecki, J.~Perolat, M.~Jaderberg and T.~Graepel.
Open-ended learning in symmetric zero-sum games. In \emph{ICML}, 2019.
\href{https://arxiv.org/abs/1901.08106}{arXiv:1901.08106}.

\bibitem[Blume and Easley(1992)]{blume1992evolution}
L.~Blume and D.~Easley. Evolution and market behavior. \emph{Journal of Economic Theory},
58:9--40, 1992. \href{https://doi.org/10.1016/0022-0531(92)90099-4}{doi:10.1016/0022-0531(92)90099-4}.

\bibitem[Blume and Easley(2006)]{blume2006smart}
L.~Blume and D.~Easley. If you're so smart, why aren't you rich? Belief selection in complete
and incomplete markets. \emph{Econometrica}, 74:929--966, 2006.
\href{https://doi.org/10.1111/j.1468-0262.2006.00691.x}{doi:10.1111/j.1468-0262.2006.00691.x}.

\bibitem[Bradley and Terry(1952)]{bradley1952rank}
R.~A. Bradley and M.~E. Terry. Rank analysis of incomplete block designs: I. The method of
paired comparisons. \emph{Biometrika}, 39:324--345, 1952.
\href{https://doi.org/10.1093/biomet/39.3-4.324}{doi:10.1093/biomet/39.3-4.324}.

\bibitem[Breiman(1961)]{breiman1961optimal}
L.~Breiman. Optimal gambling systems for favorable games. In \emph{Proceedings of the Fourth
Berkeley Symposium on Mathematical Statistics and Probability}, volume~1, 1961.
\href{https://projecteuclid.org/proceedings/berkeley-symposium-on-mathematical-statistics-and-probability/Proceedings-of-the-Fourth-Berkeley-Symposium-on-Mathematical-Statistics-and/Chapter/Optimal-gambling-systems-for-favorable-games/bsmsp/1200512159}{Project Euclid}.

\bibitem[Brier(1950)]{brier1950verification}
G.~W. Brier. Verification of forecasts expressed in terms of probability. \emph{Monthly Weather
Review}, 78(1):1--3, 1950.

\bibitem[Chiang et~al.(2024)]{chiang2024arena}
W.-L. Chiang, L.~Zheng, Y.~Sheng, A.~N. Angelopoulos, T.~Li et~al. Chatbot Arena: an open
platform for evaluating LLMs by human preference. 2024.
\href{https://arxiv.org/abs/2403.04132}{arXiv:2403.04132}.

\bibitem[Cotton(2022)]{cotton2022microprediction}
P.~Cotton. \emph{Microprediction: Building an Open AI Network}. MIT Press, 2022. Chapter~5,
``Micromanagers'', pp.~81--120.
\href{https://doi.org/10.7551/mitpress/13636.001.0001}{doi:10.7551/mitpress/13636.001.0001}.

\bibitem[Cotton(2026a)]{cotton2026wavefront}
P.~Cotton. Wavefront pruning in budgeted Brownian races. Working draft, 2026.
\url{https://brownianbandit.microprediction.org/papers.html}.

\bibitem[Cotton(2026b)]{cotton2026passk}
P.~Cotton. Posterior-predictive pass-at-$k$. Note, 2026.
\url{https://brownianbandit.microprediction.org/passk_posterior.pdf}.

\bibitem[Czarnecki et~al.(2020)]{czarnecki2020spinning}
W.~M. Czarnecki, G.~Gidel, B.~Tracey, K.~Tuyls, S.~Omidshafiei et~al. Real world games look
like spinning tops. 2020. \href{https://arxiv.org/abs/2004.09468}{arXiv:2004.09468}.

\bibitem[Gao et~al.(2022)]{gao2022overoptimization}
L.~Gao, J.~Schulman and J.~Hilton. Scaling laws for reward model overoptimization. 2022.
\href{https://arxiv.org/abs/2210.10760}{arXiv:2210.10760}.

\bibitem[Gen\c{c}ay(2026)]{gencay2026honest}
E.~Gen\c{c}ay. What survives honest evaluation? Leakage-safe, search-aware assessment of
LLM-driven trading strategy discovery. 2026.
\href{https://arxiv.org/abs/2608.27734}{arXiv:2608.27734}.

\bibitem[Gerstgrasser et~al.(2024)]{gerstgrasser2024collapse}
M.~Gerstgrasser, R.~Schaeffer, A.~Dey, R.~Rafailov et~al. Is model collapse inevitable?
Breaking the curse of recursion by accumulating real and synthetic data. 2024.
\href{https://arxiv.org/abs/2404.01413}{arXiv:2404.01413}.

\bibitem[Gneiting and Raftery(2007)]{gneiting2007proper}
T.~Gneiting and A.~E. Raftery. Strictly proper scoring rules, prediction, and estimation.
\emph{Journal of the American Statistical Association}, 102:359--378, 2007.
\href{https://doi.org/10.1198/016214506000001437}{doi:10.1198/016214506000001437}.

\bibitem[Gong et~al.(2026)]{gong2026evolutionary}
C.~Gong, J.~Bi, D.~Yao, X.~Liang et~al. Evolutionary safety of recursive self-improving AI:
taxonomy, risk discovery, and evaluation. 2026.
\href{https://arxiv.org/abs/2609.31186}{arXiv:2609.31186}.

\bibitem[Hamilton(1970)]{hamilton1970spite}
W.~D. Hamilton. Selfish and spiteful behaviour in an evolutionary model. \emph{Nature},
228:1218--1220, 1970. \href{https://doi.org/10.1038/2281218a0}{doi:10.1038/2281218a0}.

\bibitem[Hay et~al.(2012)]{hay2012selecting}
N.~Hay, S.~Russell, D.~Tolpin and S.~E. Shimony. Selecting computations: theory and
applications. In \emph{UAI}, 2012. \href{https://arxiv.org/abs/1207.5879}{arXiv:1207.5879}.

\bibitem[Huang et~al.(2023)]{huang2023selfcorrect}
J.~Huang, X.~Chen, S.~Mishra, H.~S. Zheng, A.~W. Yu, X.~Song and D.~Zhou. Large language models
cannot self-correct reasoning yet. 2023.
\href{https://arxiv.org/abs/2310.01798}{arXiv:2310.01798}.

\bibitem[Iacob et~al.(2026)]{iacob2026redqueen}
A.~Iacob, A.~Jovanovi\'c, W.~F. Shen, D.~Burkhardt et~al. The Red Queen G\"odel Machine:
co-evolving agents and their evaluators. 2026.
\href{https://arxiv.org/abs/2606.26294}{arXiv:2606.26294}.

\bibitem[Karwowski et~al.(2023)]{karwowski2023goodhart}
J.~Karwowski, O.~Hayman, X.~Bai, K.~Kiendlhofer, C.~Griffin et~al. Goodhart's law in
reinforcement learning. 2023. \href{https://arxiv.org/abs/2310.09144}{arXiv:2310.09144}.

\bibitem[Kelly(1956)]{kelly1956information}
J.~L. Kelly. A new interpretation of information rate. \emph{Bell System Technical Journal},
35:917--926, 1956.
\href{https://doi.org/10.1002/j.1538-7305.1956.tb03809.x}{doi:10.1002/j.1538-7305.1956.tb03809.x}.

\bibitem[Lanctot et~al.(2017)]{lanctot2017unified}
M.~Lanctot, V.~Zambaldi, A.~Gruslys, A.~Lazaridou, K.~Tuyls et~al. A unified game-theoretic
approach to multiagent reinforcement learning. In \emph{NeurIPS}, 2017.
\href{https://arxiv.org/abs/1711.00832}{arXiv:1711.00832}.

\bibitem[Lazear and Rosen(1981)]{lazear1981tournaments}
E.~P. Lazear and S.~Rosen. Rank-order tournaments as optimum labor contracts. \emph{Journal of
Political Economy}, 89:841--864, 1981.
\href{https://doi.org/10.1086/261010}{doi:10.1086/261010}.

\bibitem[Lichtendahl and Winkler(2007)]{lichtendahl2007competition}
K.~C. Lichtendahl and R.~L. Winkler. Probability elicitation, scoring rules, and competition
among forecasters. \emph{Management Science}, 53:1745--1755, 2007.
\href{https://doi.org/10.1287/mnsc.1070.0729}{doi:10.1287/mnsc.1070.0729}.

\bibitem[Manheim and Garrabrant(2018)]{manheim2018goodhart}
D.~Manheim and S.~Garrabrant. Categorizing variants of Goodhart's law. 2018.
\href{https://arxiv.org/abs/1803.04585}{arXiv:1803.04585}.

\bibitem[Maynard Smith and Price(1973)]{maynardsmith1973conflict}
J.~Maynard Smith and G.~R. Price. The logic of animal conflict. \emph{Nature}, 246:15--18, 1973.
\href{https://doi.org/10.1038/246015a0}{doi:10.1038/246015a0}.

\bibitem[Ng et~al.(1999)]{ng1999shaping}
A.~Y. Ng, D.~Harada and S.~Russell. Policy invariance under reward transformations: theory and
application to reward shaping. In \emph{ICML}, 1999.
\url{https://people.eecs.berkeley.edu/~russell/papers/icml99-shaping.pdf}.

\bibitem[Peters(2019)]{peters2019ergodicity}
O.~Peters. The ergodicity problem in economics. \emph{Nature Physics}, 15:1216--1221, 2019.
\href{https://doi.org/10.1038/s41567-019-0732-0}{doi:10.1038/s41567-019-0732-0}.

\bibitem[Price(1972)]{price1972fisher}
G.~R. Price. Fisher's `fundamental theorem' made clear. \emph{Annals of Human Genetics},
36:129--140, 1972. \href{https://doi.org/10.1111/j.1469-1809.1972.tb00764.x}{doi:10.1111/j.1469-1809.1972.tb00764.x}.

\bibitem[Russell and Wefald(1991)]{russell1991metareasoning}
S.~Russell and E.~Wefald. Principles of metareasoning. \emph{Artificial Intelligence},
49:361--395, 1991.
\href{https://doi.org/10.1016/0004-3702(91)90015-C}{doi:10.1016/0004-3702(91)90015-C}.

\bibitem[Sandroni(2000)]{sandroni2000markets}
A.~Sandroni. Do markets favor agents able to make accurate predictions? \emph{Econometrica},
68:1303--1341, 2000. \href{https://doi.org/10.1111/1468-0262.00163}{doi:10.1111/1468-0262.00163}.

\bibitem[Schaffer(1988)]{schaffer1988finite}
M.~E. Schaffer. Evolutionarily stable strategies for a finite population and a variable contest
size. \emph{Journal of Theoretical Biology}, 132:469--478, 1988.
\href{https://doi.org/10.1016/s0022-5193(88)80085-7}{doi:10.1016/s0022-5193(88)80085-7}.

\bibitem[Shumailov et~al.(2024)]{shumailov2024collapse}
I.~Shumailov, Z.~Shumaylov, Y.~Zhao, N.~Papernot et~al. AI models collapse when trained on
recursively generated data. \emph{Nature}, 631:755--759, 2024.
\href{https://doi.org/10.1038/s41586-024-07566-y}{doi:10.1038/s41586-024-07566-y}.

\bibitem[Skalse et~al.(2022)]{skalse2022hacking}
J.~Skalse, N.~H.~R. Howe, D.~Krasheninnikov and D.~Krueger. Defining and characterizing reward
hacking. 2022. \href{https://arxiv.org/abs/2209.13085}{arXiv:2209.13085}.

\bibitem[Smith and Winkler(2006)]{smith2006optimizer}
J.~E. Smith and R.~L. Winkler. The optimizer's curse: skepticism and postdecision surprise in
decision analysis. \emph{Management Science}, 52:311--322, 2006.
\href{https://doi.org/10.1287/mnsc.1050.0451}{doi:10.1287/mnsc.1050.0451}.

\bibitem[Tang et~al.(2026)]{tang2026gai}
H.~Tang, Y.~Ma, P.~Li and Y.~Yuan. Generalized agent iteration: one formal framework for
iterative policy improvement and recursive self-improvement. 2026.
\href{https://arxiv.org/abs/2609.13406}{arXiv:2609.13406}.

\bibitem[Taylor and Jonker(1978)]{taylor1978ess}
P.~D. Taylor and L.~B. Jonker. Evolutionary stable strategies and game dynamics.
\emph{Mathematical Biosciences}, 40:145--156, 1978.
\href{https://doi.org/10.1016/0025-5564(78)90077-9}{doi:10.1016/0025-5564(78)90077-9}.

\bibitem[Turtel et~al.(2026)]{turtel2026futureaslabel}
B.~Turtel, P.~Wilczewski, D.~Franklin and K.~Skothiem. Future-as-label: scalable supervision
from real-world outcomes. 2026. \href{https://arxiv.org/abs/2601.06336}{arXiv:2601.06336}.

\bibitem[Witkowski et~al.(2021)]{witkowski2021competitions}
J.~Witkowski, R.~Freeman, J.~W. Vaughan, D.~M. Pennock and A.~Krause. Incentive-compatible
forecasting competitions. 2021. \href{https://arxiv.org/abs/2101.01816}{arXiv:2101.01816}.

\bibitem[Yuan et~al.(2024)]{yuan2024selfrewarding}
W.~Yuan, R.~Y. Pang, K.~Cho, X.~Li, S.~Sukhbaatar et~al. Self-rewarding language models. In
\emph{ICML}, 2024. \href{https://arxiv.org/abs/2401.10020}{arXiv:2401.10020}.

\bibitem[Zhang et~al.(2025)]{zhang2025dgm}
J.~Zhang, S.~Hu, C.~Lu, R.~Lange and J.~Clune. Darwin G\"odel Machine: open-ended evolution of
self-improving agents. 2025. \href{https://arxiv.org/abs/2505.22954}{arXiv:2505.22954}.

\end{thebibliography}
\end{document}
